% Reviewed translation: PDF pages 32--36 / printed pages 18--22.
% The proof of Lemma 2.1 continues in the immutable approved sample at PDF page 37.

\section{Stokes 问题的函数空间}\label{sec:function-spaces-stokes}
\setcounter{equation}{-1}

\MathBookSourcePage{32}
严格分析 Stokes 问题需要涉及向量场散度和旋度的特殊函数空间. 在几乎所有文献中, 这些空间的关键性质
都源自 De Rham [68] 证明的一个强有力且困难的 Theorem, 其基本内容是:
\begin{equation}\label{eq:de-rham-distributional-principle}
\left\{
\begin{minipage}{0.78\linewidth}
如果分布向量场 $\mathbf u$ 对所有无散度的 $\mathcal D$ 函数都满足
$\langle\mathbf u,\boldsymbol\phi\rangle=0$, 那么对某个分布 $S$ 有
$\mathbf u=\mathbf{grad}\,S$.
\end{minipage}
\right.
\end{equation}
不过, 这个 Theorem 远非必需; 本节给出一种受 Tartar [78] 启发的不同方法.

\setcounter{subsection}{0}
\subsection{预备结果}\label{subsec:stokes-spaces-preliminary-results}

下面这个由 Peetre [63] 和 Tartar [78] 给出的 Theorem 将在本书中多次使用.

\setcounter{statement}{0}
\renewcommand{\theHstatement}{theorem.\thesection.\arabic{statement}}
\begin{Theorem}\label{thm:peetre-tartar-compact-operator}
设 $E_1,E_2,E_3$ 是三个 Banach 空间,
$A\in\mathcal L(E_1;E_2)$, 且 $B$ 是 $\mathcal L(E_1;E_3)$ 中的紧算子, 并满足
\begin{equation}\label{eq:peetre-tartar-base-equivalence}
\|u\|_{E_1}\cong\|Au\|_{E_2}+\|Bu\|_{E_3}\qquad\forall u\in E_1.
\end{equation}
则有如下性质.

$1^\circ)$ $\operatorname{Ker}(A)$ 的维数有限; 映射 $A$ 是从
$E_1/\operatorname{Ker}(A)$ 到 $\mathcal R(A)$ 的同构; $\mathcal R(A)$ 是 $E_2$ 的闭子空间.

$2^\circ)$ 存在常数 $C_0$, 使得如果 $F$ 是 Banach 空间且
$L\in\mathcal L(E_1;F)$ 在 $\operatorname{Ker}(A)$ 上为零, 则
\begin{equation}\label{eq:peetre-tartar-kernel-vanishing-bound}
\|Lu\|_F\leq C_0\|L\|_{\mathcal L(E_1;F)}\|Au\|_{E_2}\qquad\forall u\in E_1.
\end{equation}

$3^\circ)$ 如果 $G$ 是 Banach 空间且 $M\in\mathcal L(E_1;G)$ 满足
\begin{equation}\label{eq:peetre-tartar-kernel-separation}
Mu\neq0\qquad\forall u\in\operatorname{Ker}(A)-\{0\},
\end{equation}
则
\begin{equation}\label{eq:peetre-tartar-augmented-equivalence}
\|u\|_{E_1}\cong\|Au\|_{E_2}+\|Mu\|_G.
\end{equation}
\end{Theorem}

\begin{Proof}
$1^\circ)$ 这里使用一个熟知结果(参见 Taylor [80]): 如果赋范线性空间 $V$ 的单位球面是紧的,
则 $V$ 是有限维的. 将该结果应用于 $\operatorname{Ker}(A)$. 首先注意到, 由
\eqref{eq:peetre-tartar-base-equivalence}, $\|u\|_{E_1}$ 和 $\|Bu\|_{E_3}$ 是
$\operatorname{Ker}(A)$ 上的两个等价范数. 现在设 $(u_n)$ 是 $\operatorname{Ker}(A)$ 中的有界序列.
由于 $B$ 紧, 可以抽取子序列 $(u_\mu)$, 使得 $Bu_\mu$ 在 $E_3$ 中收敛. 因此 $(u_\mu)$ 是
$E_1$ 中的 Cauchy 序列, 从而是 $\operatorname{Ker}(A)$ 中的收敛序列. 所以
$\operatorname{Ker}(A)$ 的单位球面是紧的, 因而 $\operatorname{Ker}(A)$ 的维数有限.

\MathBookSourcePage{33}
既然 $\operatorname{Ker}(A)$ 是 $E_1$ 的有限维子空间, 就可以引入商空间
\[
X=E_1/\operatorname{Ker}(A),
\]
它关于通常的商范数
\[
\|\dot u\|_X=\inf_{u\in\dot u}\|u\|_{E_1}
\]
是 Banach 空间. 此外, 因为 $\operatorname{Ker}(A)$ 有限维, 上述下确界可以达到, 即每个等价类
$\dot u$ 都有一个代表元 $\tilde u$, 使得
\[
\|\dot u\|_X=\|\tilde u\|_{E_1}\qquad\forall\dot u\in X.
\]
由于 $A$ 是从 $X$ 到 $\mathcal R(A)$ 的线性连续一一映射, 为建立所述同构, 只需证明 $A$ 的逆连续,
即存在常数 $C>0$, 使得
\begin{equation}\label{eq:peetre-tartar-quotient-inverse-bound}
\|\dot u\|_X\leq C\|A\dot u\|_{E_2}\qquad\forall\dot u\in X.
\end{equation}
用反证法证明. 假设 $X$ 中存在序列 $(\dot u_n)$, 满足
\[
\|\dot u_n\|_X=1,\qquad\lim_{n\to\infty}A\dot u_n=0\quad\text{in }E_2.
\]
于是
\[
\|\tilde u_n\|_{E_1}=1,\qquad
\lim_{n\to\infty}\|A\tilde u_n\|_{E_2}=0.
\]
因此可以抽取子序列 $(\tilde u_\mu)$, 使得 $(B\tilde u_\mu)$ 在 $E_3$ 中收敛.
由 \eqref{eq:peetre-tartar-base-equivalence}, $(\tilde u_\mu)$ 是 $E_1$ 中的 Cauchy 序列.
所以存在 $u\in\operatorname{Ker}(A)$, 使得 $\tilde u_\mu\to u$ in $E_1$.
这意味着 $\dot u_\mu\to0$, 与假设矛盾.

最后, 因为 $A$ 是从 $X$ 到 $\mathcal R(A)$ 的同构且 $X$ 是 Banach 空间, 立即得到
$\mathcal R(A)$ 也是 Banach 空间. 因此 $\mathcal R(A)$ 是 $E_2$ 的闭子空间.

$2^\circ)$ 因为 $L$ 在 $\operatorname{Ker}(A)$ 上为零, 可以写成
\[
Lu=L\dot u=LA^{-1}Au\qquad\forall u\in E_1.
\]
因此
\[
\|Lu\|_F\leq\|L\|_{\mathcal L(E_1;F)}
\|A^{-1}\|_{\mathcal L(\mathcal R(A);X)}\|Au\|_{E_2}\qquad\forall u\in E_1,
\]
从而得到 \eqref{eq:peetre-tartar-kernel-vanishing-bound}, 其中
$C_0=\|A^{-1}\|_{\mathcal L(\mathcal R(A);X)}$.

$3^\circ)$ 最后证明存在常数 $C>0$, 使得
\[
\|Au\|_{E_2}+\|Mu\|_G\geq C\|u\|_{E_1}\qquad\forall u\in E_1.
\]
仍然使用反证法. 设 $(u_n)$ 是 $E_1$ 中的序列, 满足
\[
\lim_{n\to\infty}(\|Au_n\|_{E_2}+\|Mu_n\|_G)=0,\qquad\|u_n\|_{E_1}=1.
\]

\MathBookSourcePage{34}
于是存在子序列 $(u_\mu)$, 使得 $(Bu_\mu)$ 在 $E_3$ 中收敛. 因此 $(u_\mu)$ 是 $E_1$ 中的
Cauchy 序列, 并且
\[
\lim_{\mu\to\infty}u_\mu=u,
\]
其中
\[
Au=0,\qquad Mu=0,\qquad\|u\|_{E_1}=1.
\]
但是 \eqref{eq:peetre-tartar-kernel-separation} 蕴含 $u=0$, 这导致矛盾.
\end{Proof}

现在设 $\Omega$ 是 $\mathbb R^N$ 的有界子集, 其边界 $\Gamma$ 是 Lipschitz 连续的.
从现在起经常处理向量值函数. 用粗体字符区分向量, 并把前面的所有范数自然地推广到向量:
如果 $\mathbf v=(v_1,\ldots,v_N)$, 则
\[
\|\mathbf v\|_{m,p,\Omega}=\left(\sum_{i=1}^N\|v_i\|_{m,p,\Omega}^p\right)^{1/p}.
\]

下一个 Theorem 是 Nečas [57] 的一般泛函分析结果的一部分. 它的证明既长又精细, 因为只假设边界
是 Lipschitz 连续的. 当边界光滑时, 可以在 Duvaut \& Lions [26] 中找到另一种更简单的证明.
这里略去这两种证明, 因为它们都超出本书范围.

\setcounter{statement}{1}
\renewcommand{\theHstatement}{theorem.\thesection.\arabic{statement}}
\begin{Theorem}\label{thm:necas-gradient-negative-norm}
设 $\Omega$ 是有界 Lipschitz 连续开集. 存在只依赖于 $\Omega$ 的常数 $C>0$, 使得
\begin{equation}\label{eq:necas-gradient-negative-norm}
\|p\|_{0,\Omega}\leq C\{\|p\|_{-1,\Omega}+\|\mathbf{grad}\,p\|_{-1,\Omega}\}
\qquad\forall p\in L^2(\Omega).
\end{equation}
\end{Theorem}

立即得到如下结论.

\setcounter{statement}{0}
\renewcommand{\theHstatement}{corollary.\thesection.\arabic{statement}}
\begin{Corollary}\label{cor:gradient-closed-range-quotient-local}
$1^\circ)$ 在 Theorem~\ref{thm:necas-gradient-negative-norm} 的假设下, 梯度算子
\[
\mathbf{grad}\in\mathcal L(L^2(\Omega);H^{-1}(\Omega)^N)
\]
的值域是 $H^{-1}(\Omega)^N$ 的闭子空间.

$2^\circ)$ 如果进一步假设 $\Omega$ 连通, 则存在只依赖于 $\Omega$ 的常数 $C>0$, 使得
\begin{equation}\label{eq:gradient-quotient-negative-norm}
\|\dot p\|_{L^2(\Omega)/\mathbb R}\leq C\|\mathbf{grad}\,\dot p\|_{-1,\Omega}
\qquad\forall\dot p\in L^2(\Omega)/\mathbb R.
\end{equation}

$3^\circ)$ 设 $\omega$ 是 $\Omega$ 中具有正测度的开子集. 存在只依赖于 $\Omega$ 和 $\omega$ 的常数
$C_\omega>0$, 使得
\begin{equation}\label{eq:gradient-local-control}
\|p\|_{0,\Omega}\leq C_\omega\{\|p\|_{0,\omega}+\|\mathbf{grad}\,p\|_{-1,\Omega}\}
\qquad\forall p\in L^2(\Omega).
\end{equation}
\end{Corollary}

\begin{Proof}
其思想是应用 Theorem~\ref{thm:peetre-tartar-compact-operator}, 取
$E_1=L^2(\Omega)$, $E_2=H^{-1}(\Omega)^N$, $E_3=H^{-1}(\Omega)$,
$A=\mathbf{grad}$, 并令 $B$ 为 $L^2(\Omega)$ 到 $H^{-1}(\Omega)$ 的典范嵌入;
由 Theorem~\ref{thm:sobolev-embedding}, 该嵌入是紧的. 显然,
\[
\|p\|_{-1,\Omega}+\|\mathbf{grad}\,p\|_{-1,\Omega}
\leq C_1\|p\|_{0,\Omega}\qquad\forall p\in L^2(\Omega).
\]

\MathBookSourcePage{35}
因此, 由 \eqref{eq:necas-gradient-negative-norm} 立即得到
\eqref{eq:peetre-tartar-base-equivalence}, 从而证明第一个结论.

类似地, 由于 $\Omega$ 连通, $\mathbb R=\operatorname{Ker}(\mathbf{grad})$, 而
\eqref{eq:peetre-tartar-quotient-inverse-bound} 证明 \eqref{eq:gradient-quotient-negative-norm}.

最后, 令 $G=L^2(\omega)$, 并令 $M:L^2(\Omega)\to L^2(\omega)$ 为恒等映射.
因为 $\omega$ 具有正测度, 对所有非零常数 $c$ 都有 $Mc\neq0$, 而
\eqref{eq:peetre-tartar-augmented-equivalence} 蕴含 \eqref{eq:gradient-local-control}.
\end{Proof}

第二个 Corollary 给出一个重要的正则性结果.

\setcounter{statement}{1}
\renewcommand{\theHstatement}{corollary.\thesection.\arabic{statement}}
\begin{Corollary}\label{cor:local-l2-gradient-global-l2}
设 $\Omega$ 连通并满足 Theorem~\ref{thm:necas-gradient-negative-norm} 的假设. 如果
\[
p\in L^2_{\mathrm{loc}}(\Omega),\qquad\mathbf{grad}\,p\in H^{-1}(\Omega)^N,
\]
则 $p\in L^2(\Omega)$.
\end{Corollary}

\begin{Proof}
令
\[
X=\{p\in L^2_{\mathrm{loc}}(\Omega);\ \mathbf{grad}\,p\in H^{-1}(\Omega)^N\},
\]
并置
\[
\llbracket p\rrbracket_\omega=\|p\|_{0,\omega}+\|\mathbf{grad}\,p\|_{-1,\Omega},
\]
其中 $\omega\Subset\Omega$ 具有正测度. 则 $\llbracket\cdot\rrbracket_\omega$ 是 $X$ 上的范数,
由 \eqref{eq:gradient-local-control} 可知, $\llbracket\cdot\rrbracket_\omega$ 和
$\|\cdot\|_{0,\Omega}$ 是 $L^2(\Omega)$ 上的两个等价范数. 因此 $L^2(\Omega)$ 关于范数
$\llbracket\cdot\rrbracket_\omega$ 是 Banach 空间, 只需证明 $L^2(\Omega)$ 在 $X$ 中关于该范数稠密.

$1^\circ)$ 暂且假设 $\Omega$ 关于其中一点 $y$ 严格星形. 取 $y$ 为原点, 这等价于
\[
\theta\bar\Omega\subset\Omega\quad\forall\theta\in[0,1),
\qquad\bar\Omega\subset\theta\Omega\quad\forall\theta>1.
\]
取 $\theta>1$ 并令 $\Omega_\theta=\theta\Omega$. 对 $\Omega$ 上的连续函数 $\phi$,
作变量变换 $\phi\mapsto\phi_\theta$, 其中
\[
\phi_\theta(x)=\phi(x/\theta)\qquad\forall x\in\Omega_\theta.
\]
把它推广到分布 $u\in\mathcal D'(\Omega)\mapsto u_\theta\in\mathcal D'(\Omega_\theta)$:
\[
\langle u_\theta,\phi\rangle=\theta^N\langle u,\phi_{1/\theta}\rangle
\qquad\forall\phi\in\mathcal D(\Omega_\theta).
\]
容易验证
\[
\mathbf{grad}(u_\theta)=(1/\theta)(\mathbf{grad}\,u)_\theta
\qquad\forall u\in\mathcal D'(\Omega),
\]
\[
\lim_{\theta\to1}\|u_\theta-u\|_{0,\Omega}=0\qquad\forall u\in L^2(\Omega),
\]
\[
\lim_{\theta\to1}\|u_\theta-u\|_{-1,\Omega}=0\qquad\forall u\in H^{-1}(\Omega).
\]
因此, 如果 $p\in X$, 则对所有 $\theta>1$ 都有 $p_\theta\in L^2(\Omega)$, 并由上述结论立即得到

\MathBookSourcePage{36}
\[
\lim_{\theta\to1}\llbracket p_\theta-p\rrbracket_\omega=0.
\]
所以 $p\in L^2(\Omega)$.

$2^\circ)$ 在一般情形下, 使用如下性质(例如参见 Bernardi [8]):

\emph{有界 Lipschitz 连续开集是有限多个星形 Lipschitz 连续开集的并.}

显然, 只需对这些集合逐一应用上述论证, 即可在整个定义域上得到所需结果.
\end{Proof}

\setcounter{subsection}{1}
\subsection{与散度算子有关的空间的若干性质}\label{subsec:sec2-divergence-spaces}

除非另有说明, 本节假设 $\Omega$ 是 $\mathbb R^N$ 的有界子集, 其边界 $\Gamma$ 是 Lipschitz 连续的.

对 $\mathbf v=(v_1,\ldots,v_N)$, 定义散度算子
\[
\operatorname{div}\mathbf v=\sum_{i=1}^N(\partial v_i/\partial x_i).
\]
注意恒等式
\[
\operatorname{div}(\mathbf{grad}\,v)=\Delta v.
\]
引入如下无散度函数空间:
\[
\mathcal V=\{\boldsymbol\phi\in\mathcal D(\Omega)^N;\ \operatorname{div}\boldsymbol\phi=0\},
\qquad
V=\{\mathbf v\in H_0^1(\Omega)^N;\ \operatorname{div}\mathbf v=0\}.
\]
这里为 $H_0^1(\Omega)^N$ 赋予范数 $|\cdot|_{1,\Omega}$; 由
Poincaré Theorem~\ref{thm:poincare-friedrichs}, 它与 $\|\cdot\|_{1,\Omega}$ 等价.
由于 $V$ 是 $H_0^1(\Omega)^N$ 的闭子空间, 有分解
\[
H_0^1(\Omega)^N=V\oplus V^\perp,
\]
其中 $V^\perp$ 表示 $H_0^1(\Omega)^N$ 中关于与 $|\cdot|_{1,\Omega}$ 相联系的标量积
$(\mathbf{grad}\,\mathbf u,\mathbf{grad}\,\mathbf v)$ 的 $V$ 的正交补.

下述 Lemma 建立 De Rham Theorem~\eqref{eq:de-rham-distributional-principle} 的第一个粗略版本.

\setcounter{statement}{0}
\renewcommand{\theHstatement}{lemma.\thesection.\arabic{statement}}
\begin{Lemma}\label{lem:sec2-hminus1-gradient-coarse}
如果 $\mathbf f\in H^{-1}(\Omega)^N$ 满足
\begin{equation}\label{eq:sec2-annihilates-v}
\langle\mathbf f,\mathbf v\rangle=0\qquad\forall\mathbf v\in V,
\end{equation}
则存在 $p\in L^2(\Omega)$, 使得
\[
\mathbf f=\mathbf{grad}\,p.
\]
当 $\Omega$ 连通时, $p$ 在相差一个加法常数的意义下唯一.
\end{Lemma}
